\documentclass[12pt, a4paper]{ctexart}
\usepackage{amsmath, amssymb, amsthm}
\usepackage{geometry}
\usepackage{fancyhdr}
\usepackage{enumitem}
\usepackage{tikz}
\usepackage{pgfplots}
\usepackage{array}
\usepackage{multirow}
\usepackage{booktabs}

\geometry{left=2cm, right=2cm, top=2.5cm, bottom=2.5cm}

\setCJKmainfont{Hiragino Sans GB}
\setCJKsansfont{Heiti SC}

\pagestyle{fancy}
\fancyhf{}
\fancyhead[L]{圆锥曲线的方程 测试卷}
\fancyhead[R]{数学（选择性必修第一册）}
\fancyfoot[C]{第 \thepage\ 页}
\renewcommand{\headrulewidth}{0.4pt}

\setlength{\parindent}{0pt}

\newcommand{\score}[1]{\hfill\textbf{（#1分）}}

\begin{document}

\begin{center}
{\Large\bfseries 圆锥曲线的方程 测试卷}\\[0.5em]
{\normalsize 时间：120分钟 \quad 满分：150分}
\end{center}

\vspace{1em}

\noindent\textbf{一、单项选择题（本大题共8小题，每小题5分，共40分。每小题只有一个选项符合题目要求。）}

\vspace{0.5em}

\begin{enumerate}
\item 椭圆 $\dfrac{x^2}{25}+\dfrac{y^2}{9}=1$ 的离心率为 \score{5}
\begin{enumerate}[label=\Alph*.]
\item $\dfrac{3}{5}$
\item $\dfrac{4}{5}$
\item $\dfrac{5}{4}$
\item $\dfrac{5}{3}$
\end{enumerate}

\item 双曲线 $\dfrac{x^2}{4}-\dfrac{y^2}{3}=1$ 的渐近线方程为 \score{5}
\begin{enumerate}[label=\Alph*.]
\item $y=\pm\dfrac{\sqrt{3}}{2}x$
\item $y=\pm\dfrac{2}{\sqrt{3}}x$
\item $y=\pm\dfrac{3}{4}x$
\item $y=\pm\dfrac{4}{3}x$
\end{enumerate}

\item 抛物线 $y^2=8x$ 的焦点坐标为 \score{5}
\begin{enumerate}[label=\Alph*.]
\item $(2,0)$
\item $(4,0)$
\item $(0,2)$
\item $(0,4)$
\end{enumerate}

\item 椭圆 $\dfrac{x^2}{16}+\dfrac{y^2}{12}=1$ 的焦点为 $F_1$，$F_2$，$P$ 为椭圆上一点，若 $|PF_1|=5$，则 $|PF_2|=$ \score{5}
\begin{enumerate}[label=\Alph*.]
\item $3$
\item $5$
\item $7$
\item $8$
\end{enumerate}

\item 双曲线 $\dfrac{x^2}{9}-\dfrac{y^2}{16}=1$ 的焦距为 \score{5}
\begin{enumerate}[label=\Alph*.]
\item $5$
\item $8$
\item $10$
\item $\sqrt{7}$
\end{enumerate}

\item 抛物线 $y^2=4x$ 上的点到焦点的距离为 $3$，则该点的横坐标为 \score{5}
\begin{enumerate}[label=\Alph*.]
\item $1$
\item $2$
\item $3$
\item $4$
\end{enumerate}

\item 椭圆 $\dfrac{x^2}{a^2}+\dfrac{y^2}{b^2}=1$（$a>b>0$）的左顶点为 $A$，上顶点为 $B$，若 $\triangle AOB$ 为等腰直角三角形（$O$ 为原点），则椭圆的离心率为 \score{5}
\begin{enumerate}[label=\Alph*.]
\item $\dfrac{\sqrt{2}}{2}$
\item $\dfrac{\sqrt{3}}{2}$
\item $\dfrac{1}{2}$
\item $\dfrac{\sqrt{5}}{3}$
\end{enumerate}

\item 已知 $F_1$，$F_2$ 是双曲线 $\dfrac{x^2}{a^2}-\dfrac{y^2}{b^2}=1$（$a>0$，$b>0$）的两个焦点，$P$ 为双曲线上一点，$\angle F_1PF_2=60°$，$\triangle F_1PF_2$ 的面积为 $3\sqrt{3}$，则 $b=$ \score{5}
\begin{enumerate}[label=\Alph*.]
\item $\sqrt{3}$
\item $2$
\item $2\sqrt{3}$
\item $3$
\end{enumerate}
\end{enumerate}

\vspace{1em}

\noindent\textbf{二、多项选择题（本大题共4小题，每小题5分，共20分。每小题有多个选项符合题目要求，全部选对得5分，部分选对得2分，有选错的得0分。）}

\vspace{0.5em}

\begin{enumerate}[resume]
\item 关于椭圆 $\dfrac{x^2}{4}+y^2=1$，下列结论正确的是 \score{5}
\begin{enumerate}[label=\Alph*.]
\item 长轴长为 $4$
\item 短轴长为 $2$
\item 离心率为 $\dfrac{\sqrt{3}}{2}$
\item 焦距为 $2\sqrt{3}$
\end{enumerate}

\item 关于双曲线 $\dfrac{x^2}{4}-\dfrac{y^2}{5}=1$，下列结论正确的是 \score{5}
\begin{enumerate}[label=\Alph*.]
\item 实轴长为 $4$
\item 虚轴长为 $2\sqrt{5}$
\item 离心率为 $\dfrac{3}{2}$
\item 焦点到渐近线的距离为 $\sqrt{5}$
\end{enumerate}

\item 关于抛物线 $y^2=8x$，下列结论正确的是 \score{5}
\begin{enumerate}[label=\Alph*.]
\item 焦点为 $(2,0)$
\item 准线方程为 $x=-2$
\item 通径长为 $8$
\item 离心率为 $1$
\end{enumerate}

\item 已知椭圆 $\dfrac{x^2}{16}+\dfrac{y^2}{12}=1$，点 $A(-2,0)$，$B(2,0)$，$P$ 为椭圆上动点，则 \score{5}
\begin{enumerate}[label=\Alph*.]
\item $|PA|+|PB|=8$
\item $|PA|\cdot|PB|$ 的最大值为 $16$
\item $\overrightarrow{PA}\cdot\overrightarrow{PB}$ 的最小值为 $8$
\item $\triangle PAB$ 面积的最大值为 $6\sqrt{3}$
\end{enumerate}
\end{enumerate}

\vspace{1em}

\noindent\textbf{三、填空题（本大题共4小题，每小题5分，共20分。）}

\vspace{0.5em}

\begin{enumerate}[resume]
\item 椭圆 $\dfrac{x^2}{25}+\dfrac{y^2}{16}=1$ 的长轴长为 \underline{\hspace{2cm}}，短轴长为 \underline{\hspace{2cm}}，离心率为 \underline{\hspace{2cm}}。\score{5}

\item 双曲线 $\dfrac{x^2}{9}-\dfrac{y^2}{4}=1$ 的实轴长为 \underline{\hspace{2cm}}，虚轴长为 \underline{\hspace{2cm}}，渐近线方程为 \underline{\hspace{3cm}}。\score{5}

\item 过抛物线 $y^2=12x$ 的焦点且倾斜角为 $45°$ 的弦长为 \underline{\hspace{2cm}}。\score{5}

\item 已知椭圆 $\dfrac{x^2}{a^2}+\dfrac{y^2}{b^2}=1$（$a>b>0$）的离心率为 $\dfrac{1}{2}$，且过点 $(2,\sqrt{3})$，则椭圆方程为 \underline{\hspace{5cm}}。\score{5}
\end{enumerate}

\vspace{1em}

\noindent\textbf{四、解答题（本大题共6小题，共70分。解答应写出文字说明、证明过程或演算步骤。）}

\vspace{0.5em}

\begin{enumerate}[resume]
\item （10分）已知椭圆的中心在原点，焦点在 $x$ 轴上，长轴长为 $10$，离心率为 $\dfrac{3}{5}$。
\begin{enumerate}[label=(\arabic*)]
\item 求椭圆的标准方程；
\item 求椭圆上到点 $(3,0)$ 距离最近的点的坐标。
\end{enumerate}

\item （10分）已知双曲线 $\dfrac{x^2}{4}-\dfrac{y^2}{12}=1$。
\begin{enumerate}[label=(\arabic*)]
\item 求双曲线的焦点坐标、离心率和渐近线方程；
\item 求双曲线上的点到焦点 $(4,0)$ 距离的最小值。
\end{enumerate}

\item （12分）已知抛物线 $y^2=2px$（$p>0$）的焦点为 $F$，点 $M(2,y_0)$ 在抛物线上，且 $|MF|=3$。
\begin{enumerate}[label=(\arabic*)]
\item 求抛物线的方程；
\item 过焦点 $F$ 的直线 $l$ 与抛物线交于 $A$，$B$ 两点，若 $|AB|=10$，求直线 $l$ 的方程。
\end{enumerate}

\item （12分）已知椭圆 $\dfrac{x^2}{a^2}+\dfrac{y^2}{b^2}=1$（$a>b>0$）经过点 $(0,\sqrt{3})$，且离心率为 $\dfrac{1}{2}$。
\begin{enumerate}[label=(\arabic*)]
\item 求椭圆方程；
\item 设直线 $l: y=kx+1$ 与椭圆交于 $A$，$B$ 两点，求 $\triangle AOB$ 面积的最大值（$O$ 为原点）。
\end{enumerate}

\item （12分）已知点 $A(0,2)$，$B(0,-2)$，动点 $P$ 满足 $|PA|-|PB|=2\sqrt{3}$。
\begin{enumerate}[label=(\arabic*)]
\item 求动点 $P$ 的轨迹方程；
\item 过点 $B$ 作斜率为 $k$ 的直线与轨迹交于 $C$，$D$ 两点，求 $\triangle ACD$ 面积的最小值。
\end{enumerate}

\item （14分）已知椭圆 $\dfrac{x^2}{6}+\dfrac{y^2}{2}=1$，$F_1$，$F_2$ 分别为左、右焦点。
\begin{enumerate}[label=(\arabic*)]
\item 设椭圆上一点 $P$ 满足 $\angle F_1PF_2=90°$，求 $\triangle F_1PF_2$ 的面积；
\item 过点 $F_2$ 的直线 $l$ 与椭圆交于 $A$，$B$ 两点，若以 $AB$ 为直径的圆过点 $F_1$，求直线 $l$ 的方程；
\item 在（2）的条件下，求 $\triangle ABF_1$ 的面积。
\end{enumerate}
\end{enumerate}

\end{document}
